\cgeentry{cge_nest}{Declare a nested CES or CET block}

\cgesec{Description}

\cmd{cge_nest} (from version 0.11.0) declares a production or transformation
function as a tree of CES or CET nests over variables of the model, as the
production blocks of MPSGE do \citep{rutherford1999}. \cmd{cge_nest begin}
names the block, its output quantity and price and the elasticity at the top:
\opt{}{s()} for a cost block, whose inputs substitute for each other, or
\opt{}{t()} for a revenue block, whose outputs are transformed into each
other. \cmd{cge_nest nest} adds a nest within the block or within another
nest. \cmd{cge_nest input} (in a cost block) and \cmd{cge_nest output} (in a
revenue block) add a leaf: a variable with the expression of its price.
\cmd{cge_nest end} writes the block as user variables and equations
(\manref{cge_variable}, \manref{cge_equation}) in calibrated share form, with
the shares and base values from \cmd{base()}, so that the block holds at the
base and follows the model when it is calibrated again. The leaves' equations
and the price of the output take the place of the built-in equations named in
\opt{rep}{lace()}.

\cgesec{Quick start}

Value added of labour and capital at an elasticity of 0.8, in place of the
built-in CES function (Q3) and factor demands (Q4)
\begin{stcode}
cge_nest begin va(a), output(QVA(a)) price(PVA(a)) s(0.8) replace(Q3)
cge_nest input QF(lab, a), price(W(lab) * WFDIST(lab, a)) replace(Q4(lab, a))
cge_nest input QF(cap, a), price(W(cap) * WFDIST(cap, a)) replace(Q4(cap, a))
cge_nest end
\end{stcode}
Labour in a nest of its own within the block
\begin{stcode}
cge_nest nest labour, s(1.5)
cge_nest input QF(lab_hi, a), price(W(lab_hi) * WFDIST(lab_hi, a)) in(labour) replace(Q4(lab_hi, a))
\end{stcode}
List the blocks; drop the block that is open
\begin{stcode}
cge_nest, list
cge_nest, drop
\end{stcode}

\cgesec{Syntax}

Begin a block, add nests and leaves, end it
\begin{stsyntax}
\cmdb{cge_nest begin} \argu{name}[\cmd{(}\argu{i}\cmd{)}][\cmd{\$(}\argu{condition}\cmd{)}]\texttt{,}
\optarg{out}{put}{reference} \optarg{pr}{ice}{reference} \{\optarg{}{s}{\#}|\optarg{}{t}{\#}\} [\optarg{eff}{iciency}{expression} \optarg{rep}{lace}{block}]\\
\cmdb{cge_nest nest} \argu{nestname}\texttt{,} \{\optarg{}{s}{\#}|\optarg{}{t}{\#}\} [\optarg{}{in}{nestname}]\\
\cmdb{cge_nest input} \argu{reference}\texttt{,} \optarg{pr}{ice}{expression} [\optarg{un}{its}{expression} \optarg{}{in}{nestname} \optarg{rep}{lace}{block}]\\
\cmdb{cge_nest output} \argu{reference}\texttt{,} \optarg{pr}{ice}{expression} [\optarg{un}{its}{expression} \optarg{}{in}{nestname} \optarg{rep}{lace}{block}]\\
\cmdb{cge_nest end}
\end{stsyntax}
List the blocks, or drop the block that is open
\begin{stsyntax}
\cmdb{cge_nest} \texttt{,} \{\opt{}{list}|\opt{}{drop}\}
\end{stsyntax}
\argu{name} has at most 12 characters; the index \argu{i} and the condition
are those of \manref{cge_equation}. A \argu{reference} is a variable,
\cmd{QF(lab, a)}, with the block's index; an \argu{expression} is one of the
language of \manref{cge_equation}.

\begin{opttable}
\optgroup{begin}
\optarg{out}{put}{reference} & the output quantity\\
\optarg{pr}{ice}{reference} & the price of the output\\
\optarg{}{s}{\#}, \optarg{}{t}{\#} & elasticity of substitution (cost block) or of transformation (revenue block)\\
\optarg{eff}{iciency}{expression} & efficiency, 1 at the base; default \cmd{1}\\
\optarg{rep}{lace}{block} & the built-in equation that the price of the output replaces\\
\optgroup{nest}
\optarg{}{s}{\#}, \optarg{}{t}{\#} & elasticity of the nest\\
\optarg{}{in}{nestname} & the nest it is within; default the block\\
\optgroup{input, output}
\optarg{pr}{ice}{expression} & the price of the leaf\\
\optarg{un}{its}{expression} & the units in which the leaf enters; default \cmd{1}\\
\optarg{}{in}{nestname} & the nest it is within; default the block\\
\optarg{rep}{lace}{block} & the built-in equation that the leaf's equation replaces\\
\end{opttable}

\cgesec{Options}

\begin{optdesc}
\item[\optarg{out}{put}{reference} \optarg{pr}{ice}{reference}] name the
variables of the block's output and its price, with the block's index, as
\cmd{output(QVA(a)) price(PVA(a))}.

\item[\optarg{}{s}{\#}, \optarg{}{t}{\#}] give the elasticity, 0 or more, of
substitution in a cost block and its nests, or of transformation in a revenue
block and its nests. At $s = 1$ a nest is Cobb-Douglas, at 0 Leontief.

\item[\optarg{eff}{iciency}{expression}] scales the output of the block; it
should be 1 at the base, as \cmd{efficiency(ad(a) / base(ad(a)))}, so that a
shock to \cmd{ad} moves the block as it moves the built-in value added.

\item[\optarg{rep}{lace}{block}] names the built-in equation that the
equation is put in place of, as \opt{rep}{lace()} of \manref{cge_equation}.
With the block's index, a label may stand beside it, as
\cmd{replace(Q4(lab, a))}: the element of Q4 for labour and each activity.

\item[\optarg{pr}{ice}{expression}] of a leaf is the price paid for, or
received for, its quantity, as \cmd{price(W(lab) * WFDIST(lab, a) * (1 + tf(lab, a)))}
with the tax on factor use.

\item[\optarg{un}{its}{expression}] makes a leaf enter in units of its own:
the block takes \argu{units} times the quantity, at the price over
\argu{units}, as \cmd{units(af(lab))} for a factor in efficiency units.

\item[\optarg{}{in}{nestname}] puts a nest or a leaf within a nest.

\item[\opt{}{list}] lists the declared blocks; \opt{}{drop} drops the block
that is open. \manref{cge_equation}\cmd{, clear} removes the blocks with the
other user variables and equations.
\end{optdesc}

\cgesec{Remarks and examples}

A block of value added replaces the built-in function of value added (Q3) and
the demands for factors (Q4): its leaves' equations give the factors, and its
top equation the price of value added. The SAM of the example has two kinds of
labour, capital, an energy commodity and a payroll tax on skilled labour:
\stlog{nest_setup}
The two kinds of labour are put in a nest with an elasticity of 1.5, below a
top of 0.8 with capital; the prices of the factors include the tax on their
use, and the factors enter in efficiency units:
\stlog{nest_declare}
Each leaf has a price per unit and a quantity in units, \cmd{va_w}\argu{k} and
\cmd{va_x}\argu{k}; the nest has a price index, 1 at the base, and a
quantity, \cmd{va_labour_P} and \cmd{va_labour_Q}, and the block a price
index, \cmd{va_P}. The base needs no iteration. Ten per cent more skilled
labour:
\stlog{nest_solve}
The same structure in the built-in nests of value added (\manref{cge_calibrate},
\cmd{nests()}) gives the same solution:
\stlog{nest_builtin}
The test suite checks blocks of one nest at the built-in elasticity, of two
nests with one elasticity, Cobb-Douglas blocks and CET blocks of exports and
domestic sales against the built-in functions, and this block against the
built-in nests, to $10^{-10}$; \manref{cge_gams} solves a block in GAMS.

\cgesec{Stored results}

\cmd{cge_nest end} stores the following in \cmd{r()}:
\begin{storedres}
\resgroup{Scalars}
\res{r(N_variables)} & user variables written\\
\res{r(N_equations)} & user equations written\\
\resgroup{Macros}
\res{r(name)} & name of the block\\
\end{storedres}
\cmd{cge_nest, list} stores \cmd{r(N_blocks)} and \cmd{r(blocks)}.

\cgesec{Methods and formulas}

In nest $n$ with elasticity of substitution $s_n$, children $j$ (leaves or
nests) with prices $p_j$, base prices $p^0_j$, base quantities $q^0_j$ and
base value shares $\theta_j = p^0_j q^0_j / \sum_k p^0_k q^0_k$,
\begin{equation*}
P_n = \Bigl[\sum_j \theta_j \bigl(p_j/p^0_j\bigr)^{1-s_n}\Bigr]^{1/(1-s_n)},
\qquad
q_j = q^0_j\,\frac{Q_n}{Q^0_n}\,\Bigl(\frac{P_n\,p^0_j}{p_j}\Bigr)^{s_n},
\end{equation*}
with $P_n = \prod_j (p_j/p^0_j)^{\theta_j}$ at $s_n = 1$ and
$P_n = \sum_j \theta_j\,p_j/p^0_j$, $q_j = q^0_j\,Q_n/Q^0_n$ at $s_n = 0$. A
nest's price index is 1 at the base and its quantity its value at base
prices, the sum of $p^0_j q^0_j$ over the leaves within it. A revenue nest
with elasticity of transformation $t_n$ has the exponent $1 + t_n$ in place of
$1 - s_n$ and $q_j = q^0_j\,(Q_n/Q^0_n)\,\bigl(p_j/(p^0_j P_n)\bigr)^{t_n}$. At
the top, $Q_n/Q^0_n$ is the output relative to its base, divided by
efficiency $E$ in a cost block and multiplied by it in a revenue block, and
the price of the output is its base value times $P_n/E$ (cost) or $P_n E$
(revenue). With units $u_j$, a leaf enters as $u_j\,x_j$ at $p_j/u_j$. The
shares and base values are written with \cmd{base()}, of variables and from
version 0.11.0 of parameters, so they are taken again from the base when the
model is calibrated again. Section~\ref{sec:usernests} gives the
equations written.

\cgesec{Reference}

\citet{rutherford1999}.

\cgesec{Also see}

\begin{alsosee}
\seeentry{cge_equation}{Add a user equation}
\seeentry{cge_variable}{Add a user variable}
\seeentry{cge_calibrate}{Calibrate the model}
\end{alsosee}
